\documentclass[11pt]{article}
\usepackage[margin=1in]{geometry}
\usepackage{enumitem}
% =============================================================================
% Preambles/header.tex — shared LaTeX/Beamer preamble for this project
%
% Use from a Beamer lecture by prepending:
%     \input{header}
% and compiling with TEXINPUTS=../Preambles:$TEXINPUTS (the /compile-latex
% skill does this automatically).
%
% PALETTE CONTRACT. Color names defined here MUST match the names in
% Quarto/theme-template.scss so Beamer and Quarto renderings use the same
% palette. The scripts/check-palette-sync.sh script enforces this.
%
% To customize: edit the HEX values below AND the matching $variable in
% Quarto/theme-template.scss, then run scripts/check-palette-sync.sh.
% =============================================================================

% -----------------------------------------------------------------------------
% Engine detection + encoding
%
% Our /compile-latex skill uses XeLaTeX, where inputenc is unnecessary and
% can produce encoding warnings. Only load inputenc under pdfTeX.
% -----------------------------------------------------------------------------
\usepackage{iftex}
\ifPDFTeX
  \usepackage[utf8]{inputenc}
\fi

\usepackage{amsmath, amssymb, amsthm}
\usepackage{booktabs}
\usepackage{graphicx}

% -----------------------------------------------------------------------------
% PALETTE — mirrors Quarto/theme-template.scss variable names.
% Load xcolor and define colors BEFORE anything that references them
% (notably hyperref, hypersetup, and the Beamer color assignments).
% -----------------------------------------------------------------------------
\usepackage{xcolor}

\definecolor{primary-blue}{HTML}{012169}      % headings, accents
\definecolor{primary-gold}{HTML}{B9975B}      % emphasis, borders
\definecolor{highlight-yellow}{HTML}{F2A900}  % markers, alerts
\definecolor{light-bg}{HTML}{E8EDF5}          % title / section backgrounds
\definecolor{jet}{HTML}{1A1A1A}               % body text

% Semantic colors (used in diagrams and annotations)
\definecolor{positive}{HTML}{15803D}          % good / observed / identified
\definecolor{negative}{HTML}{B91C1C}          % bad / problematic / bias
\definecolor{neutral}{HTML}{525252}           % reference / context

% Highlight accents (match .hi-* classes in the SCSS)
\definecolor{hi-slate}{HTML}{314F4F}
\definecolor{hi-green}{HTML}{2E8B57}
\definecolor{hi-red}{HTML}{C41E3A}

% Semantic aliases so skill templates and snippets can write
% \textcolor{accent}{...} without hard-coding "primary-blue".
\colorlet{accent}{primary-blue}
\colorlet{accent2}{primary-gold}

% -----------------------------------------------------------------------------
% Hyperref — loaded AFTER xcolor + palette so link colors resolve.
% -----------------------------------------------------------------------------
\usepackage{hyperref}
\hypersetup{colorlinks=true, linkcolor=primary-blue, urlcolor=primary-blue,
            citecolor=primary-blue, pdfborder={0 0 0}}

% -----------------------------------------------------------------------------
% Course code — set once per deck (after \input{header}, before \begin{document})
% via \coursecode{CS401}. Shown in the footer on every slide so a deck's course
% is identifiable without opening the title page. Empty by default.
% -----------------------------------------------------------------------------
\makeatletter
\newcommand{\coursecode}[1]{\gdef\@coursecodetext{#1}}
\gdef\@coursecodetext{}
\makeatother

% -----------------------------------------------------------------------------
% Beamer theme assignments (only apply under Beamer).
%
% We detect Beamer via \@ifclassloaded{beamer}, which is the documented,
% non-internal way. This must be wrapped in \makeatletter/\makeatother
% because \@ifclassloaded contains an @-catcode character.
% -----------------------------------------------------------------------------
\makeatletter
\@ifclassloaded{beamer}{%
  \usefonttheme{professionalfonts}
  \setbeamercolor{structure}{fg=primary-blue}
  \setbeamercolor{frametitle}{fg=primary-blue}
  \setbeamercolor{title}{fg=primary-blue}
  \setbeamercolor{subtitle}{fg=primary-gold}
  \setbeamercolor{author}{fg=primary-blue}
  \setbeamercolor{institute}{fg=neutral}
  \setbeamercolor{date}{fg=primary-gold}
  \setbeamercolor{itemize item}{fg=primary-blue}
  \setbeamercolor{itemize subitem}{fg=primary-gold}
  \setbeamercolor{alerted text}{fg=primary-gold}
  \setbeamercolor{block title}{fg=white, bg=primary-blue}
  \setbeamercolor{block body}{bg=light-bg}
  % Semantic block variants: block=definition/concept (blue), exampleblock=
  % worked example (green/positive), alertblock=key takeaway or caution (gold).
  % Keep to <=2 colored boxes per slide across ALL three variants (INV-7).
  \setbeamercolor{block title example}{fg=white, bg=positive}
  \setbeamercolor{block body example}{bg=light-bg}
  \setbeamercolor{block title alerted}{fg=white, bg=primary-gold}
  \setbeamercolor{block body alerted}{bg=light-bg}

  % Minimal footer: course code (left) + page X / N (right), no navigation clutter
  \setbeamertemplate{navigation symbols}{}
  \setbeamertemplate{footline}{%
    \color{neutral}\footnotesize\hspace{0.7em}\@coursecodetext\hfill
    \insertframenumber/\inserttotalframenumber\hspace{0.7em}\vspace{0.3em}}
}{}
\makeatother

% -----------------------------------------------------------------------------
% TikZ setup — libraries the snippet gallery uses
% -----------------------------------------------------------------------------
\usepackage{tikz}
\usetikzlibrary{arrows.meta, positioning, calc, decorations.pathreplacing,
                fit, shapes.geometric, backgrounds}

% Shared TikZ styles — snippets and hand-written diagrams can pick these up
% via `\begin{tikzpicture}[every node/.style={...}]` or by naming specific
% styles (e.g., `\node[dag-node] (X) at (0,0) {X};`).
\tikzset{
  % Node styles for causal DAGs and similar
  dag-node/.style = {
    draw=primary-blue, fill=light-bg, rounded corners=2pt,
    minimum width=1.2cm, minimum height=0.9cm, align=center, font=\small
  },
  decision-node/.style = {
    draw=primary-gold, fill=white, diamond, aspect=1.8,
    minimum width=1.8cm, minimum height=1.2cm, align=center, font=\small,
    inner sep=1pt
  },
  % Edge styles — observed vs counterfactual
  observed-edge/.style = {-{Latex[length=2.5mm]}, thick, draw=jet},
  counterfactual-edge/.style = {-{Latex[length=2.5mm]}, thick, draw=neutral, dashed},
  confound-edge/.style = {-{Latex[length=2.5mm]}, thick, draw=negative, dashed},
  % Data-point styles for plots
  observed-dot/.style = {fill=primary-blue, draw=primary-blue, circle, inner sep=1.2pt},
  counterfactual-dot/.style = {fill=white, draw=neutral, circle, inner sep=1.2pt}
}

% -----------------------------------------------------------------------------
% Convenience macros
% -----------------------------------------------------------------------------
\newcommand{\muted}[1]{\textcolor{neutral}{#1}}
\newcommand{\key}[1]{\textcolor{primary-gold}{\textbf{#1}}}
\newcommand{\good}[1]{\textcolor{positive}{#1}}
\newcommand{\bad}[1]{\textcolor{negative}{#1}}

% Standout frame macro: full-bleed dark background for section transitions.
% Beamer-only (uses \begin{frame}/\setbeamercolor/\usebeamerfont) -- do not
% call from an article-class file (e.g. Notes/<CODE>/*-notes.tex). \muted,
% \key, \good, \bad, and \coursecode above are plain xcolor/gdef and are
% safe to use from either class; verified when Notes/article-class support
% was added.
% Expand: \transitionslide{Your transition title}
\newcommand{\transitionslide}[1]{%
  \begingroup
  \setbeamercolor{background canvas}{bg=primary-blue}
  \begin{frame}[plain]
    \centering
    \vfill
    {\usebeamerfont{title}\color{white}\Large #1\par}
    \vfill
  \end{frame}
  \endgroup
}

% Section divider frame (Beamer-only), an alternative to \transitionslide with a
% darker two-line style: near-black (jet) background, a small white label line
% above a larger gold title, and a thin gold rule along the bottom edge.
% Expand: \sectiondivider{Section 2}{Pointers}
\newcommand{\sectiondivider}[2]{%
  \begingroup
  \setbeamercolor{background canvas}{bg=jet}
  \begin{frame}[plain]
    \vspace*{3.0cm}
    {\color{white}\ttfamily\large #1\par}
    \vspace{0.5cm}
    {\color{primary-gold}\LARGE #2\par}
    \begin{tikzpicture}[remember picture, overlay]
      \draw[primary-gold, line width=2.5pt]
        ($(current page.south west)+(0,0.1)$) -- ($(current page.south east)+(0,0.1)$);
    \end{tikzpicture}
  \end{frame}
  \endgroup
}

% -----------------------------------------------------------------------------
% Channel watermark — "Concepts That Click" (YouTube).
%
% Article-class files (CompetitiveExam Q/A sets, Notes, Assignments) get a
% light diagonal draft watermark on every page plus a centered footer naming
% the channel. Beamer decks get a subtle TikZ background watermark instead
% (the footline already carries the course code). Centralized here so every
% MCQ PDF is branded without per-file edits.
% -----------------------------------------------------------------------------
\makeatletter
\@ifclassloaded{article}{%
  \usepackage{draftwatermark}
  \SetWatermarkText{Concepts That Click}
  \SetWatermarkScale{1}
  \SetWatermarkFontSize{2cm}
  \SetWatermarkLightness{0.86}
  \SetWatermarkAngle{45}
  \usepackage{fancyhdr}
  \pagestyle{fancy}
  \fancyhf{}
  \fancyfoot[C]{\footnotesize\textcolor{neutral}{Concepts That Click~|~YouTube: \href{https://www.youtube.com/@concepts.thatclick}{@concepts.thatclick}}}
  \renewcommand{\headrulewidth}{0pt}
  \renewcommand{\footrulewidth}{0pt}
  \fancypagestyle{plain}{%
    \fancyhf{}%
    \fancyfoot[C]{\footnotesize\textcolor{neutral}{Concepts That Click~|~YouTube: \href{https://www.youtube.com/@concepts.thatclick}{@concepts.thatclick}}}%
    \renewcommand{\headrulewidth}{0pt}%
    \renewcommand{\footrulewidth}{0pt}%
  }
}{}
\@ifclassloaded{beamer}{%
  \setbeamertemplate{background}{%
    \begin{tikzpicture}[remember picture, overlay]
      \node[rotate=45, opacity=0.055, align=center]
        at (current page.center)
        {\Huge\textcolor{neutral}{Concepts That Click}};
    \end{tikzpicture}%
  }
}{}
\makeatother 
\coursecode{JKSSB}
\renewcommand{\thesection}{57.\arabic{section}}
\newtheorem{example}{Example}[section]
\newenvironment{definitionbox}[1]{%
  \par\noindent\textbf{Definition (#1).}\ \itshape
}{\par\vspace{0.3em}}
\title{Digital Platforms and E-Governance --- Detailed Lecture Notes}
\author{JKSSB Computer Awareness}
\date{\today}

\begin{document}
\maketitle

\section{What e-Governance means}
A running government delivers many services: certificates, licences,
payments, benefits, and information. \textbf{e-Governance} is the use of
Information and Communication Technology (ICT) by government to deliver
these services. Instead of standing at a physical counter, a citizen reaches
the service through an online portal. The aim is to make services
accessible, efficient, transparent, and reliable.

e-Governance is about governance, not merely about technology. The
technology is the tool; the purpose is better delivery of public services.

\begin{definitionbox}{e-Governance}
The use of Information and Communication Technology (ICT) by government to
deliver services to citizens, businesses, and other parts of government in
an accessible, efficient, transparent, and reliable way.
\end{definitionbox}

\begin{center}
\begin{tabular}{p{0.22\textwidth}p{0.68\textwidth}}
\toprule
\textbf{Term} & \textbf{Meaning in this lesson} \\
\midrule
e-Governance & ICT-based delivery of government services. \\
G2x & The interaction categories: G2C, G2G, G2B, G2E. \\
Digital India & The flagship Government of India programme launched on 1 July 2015. \\
NeGP & National e-Governance Plan, approved in May 2006. \\
e-Kranti & Pillar 5 of Digital India and the upgraded NeGP (NeGP 2.0). \\
NeGD & National e-Governance Division, created in 2009 under MeitY. \\
e-NAM & National Agriculture Market, the pan-India agricultural trading portal. \\
\bottomrule
\end{tabular}
\end{center}

\section{The G2x classification}
e-Governance interactions are named by who takes part. The ``G'' always
stands for Government, and the letter after it says who the government is
interacting with.

\begin{center}
\begin{tabular}{p{0.12\textwidth}p{0.30\textwidth}p{0.48\textwidth}}
\toprule
\textbf{Type} & \textbf{Full form} & \textbf{Who is served / examples} \\
\midrule
G2C & Government to Citizen & The general public. UMANG, DigiLocker, certificates, utility bills. \\
G2G & Government to Government & Government departments and levels. e-Office, PFMS. \\
G2B & Government to Business & Businesses and companies. GeM, GSTN. \\
G2E & Government to Employee & Government employees. An employee service book or salary portal. \\
\bottomrule
\end{tabular}
\end{center}

The classification is the highest-yield part of this topic. An item can
name a platform and ask which G2x category it belongs to, or name a category
and ask which platform fits. The safe method is to ask ``who is the
government serving?'' A citizen means G2C; another government body means
G2G; a company means G2B; an employee means G2E.

\begin{example}
The stem asks, ``GeM (Government e-Marketplace) is an example of which
category?'' GeM is used by government to buy goods and services from
businesses, so the government is serving a business. The answer is G2B, not
G2C. A student who sees only that ``government uses it'' may wrongly choose
G2C; the deciding question is who is on the other side.
\end{example}

\section{Digital India}
\textbf{Digital India} is a flagship programme of the Government of India.
It was launched on \textbf{1 July 2015} by the Prime Minister, with the
vision to transform India into a \textbf{digitally empowered society and
knowledge economy}.

The programme rests on \textbf{three vision areas}:
\begin{enumerate}[label=(\roman*)]
  \item Digital Infrastructure as a Core Utility to Every Citizen;
  \item Governance and Services on Demand;
  \item Digital Empowerment of Citizens.
\end{enumerate}

Under these vision areas the programme is built on \textbf{nine pillars}:
Broadband Highways; Universal Access to Mobile Connectivity; Public Internet
Access Programme; e-Governance: Reforming Government through Technology;
e-Kranti: Electronic Delivery of Services; Information for All; Electronics
Manufacturing; IT for Jobs; and Early Harvest Programmes. Pillar 5,
e-Kranti, is the service-delivery pillar and is treated separately below.

\begin{definitionbox}{Digital India}
A flagship Government of India programme, launched on 1 July 2015, with the
vision of transforming India into a digitally empowered society and
knowledge economy, built on three vision areas and nine pillars.
\end{definitionbox}

\section{NeGP, e-Kranti, and NeGD: three names, three things}
These three names are close and are a standing source of confusion. They are
kept apart here by asking what each one \emph{is}.

The \textbf{National e-Governance Plan (NeGP)} is the older umbrella
\emph{plan}. It was approved in \textbf{May 2006}, with the vision to make
all government services accessible to the common man in his locality,
through common service delivery outlets, at affordable cost. It is built
around \textbf{Mission Mode Projects (MMPs)} --- 27 MMPs with supporting
components.

\textbf{e-Kranti} means ``electronic delivery of services''. It is
\textbf{Pillar 5} of the Digital India programme and is also called
\textbf{NeGP 2.0}, an upgraded approach to the National e-Governance Plan.
Its vision is ``Transforming e-Governance for Transforming Governance'',
and its mission is to deliver all government services electronically through
integrated and interoperable systems.

The \textbf{National e-Governance Division (NeGD)} is not a plan at all. It
is an \emph{organisation}. It was created in \textbf{2009} by the Ministry
of Electronics and Information Technology (MeitY) as an independent business
division under the Digital India Corporation. Its role is to support MeitY
in the programme management and implementation of e-Governance projects, and
to give technical and advisory support to ministries and departments.

\begin{center}
\begin{tabular}{p{0.12\textwidth}p{0.40\textwidth}p{0.10\textwidth}p{0.22\textwidth}}
\toprule
\textbf{Name} & \textbf{What it is} & \textbf{Year} & \textbf{One word} \\
\midrule
NeGP & Umbrella e-Governance plan with 27 MMPs & 2006 & The plan \\
e-Kranti & Pillar 5 of Digital India; NeGP 2.0 & 2015 & The approach \\
NeGD & Division under MeitY supporting projects & 2009 & The body \\
\bottomrule
\end{tabular}
\end{center}

\begin{example}
An item states, ``NeGD is the National e-Governance Plan approved in 2006.''
This is wrong on two counts. The plan approved in 2006 is \textbf{NeGP}
(National e-Governance Plan). \textbf{NeGD} is the National e-Governance
\emph{Division}, an organisation created in 2009. The distractor swaps a
plan for a body.
\end{example}

\section{e-NAM and agricultural marketing}
\textbf{e-NAM} stands for \textbf{National Agriculture Market} (also written
eNAM). It is a pan-India \textbf{electronic trading portal} for agricultural
commodities, launched on \textbf{14 April 2016}. It is implemented by the
\textbf{Small Farmers Agribusiness Consortium (SFAC)} under the Ministry of
Agriculture and Farmers Welfare. The portal is available at
\textbf{enam.gov.in}.

The key idea is that e-NAM \textbf{networks the existing Agriculture Produce
Market Committee (APMC) mandis} to create a \textbf{unified national market}
for agricultural commodities. Because trade happens online, farmers get
\textbf{online competitive bidding and transparent price discovery}, and
direct sale to buyers reduces the role of intermediaries.

\begin{definitionbox}{e-NAM}
The National Agriculture Market, a pan-India electronic trading portal,
launched on 14 April 2016, that networks the existing APMC mandis to create
a unified national market for agricultural commodities.
\end{definitionbox}

A near-neighbour to keep separate is \textbf{AGMARKNET}. AGMARKNET is a
portal that provides agricultural marketing \emph{information} such as
market prices. e-NAM, by contrast, is the \emph{trading} platform where the
sale itself takes place. An item that asks for the agricultural marketing
\emph{trading} platform wants e-NAM.

\section{Citizen portal, electronic records, and authentication}
The \textbf{National Portal of India} is \textbf{india.gov.in}. It is a
\textbf{single-window} gateway from which a citizen can reach the
information and services of the Government of India. A related platform,
\textbf{MyGov}, is the citizen-engagement platform for participative
governance.

Electronic governance depends on giving electronic documents the same
standing as paper. The \textbf{Information Technology Act, 2000} gives legal
recognition to \textbf{electronic records}. An electronic record can be
stored, retrieved, and used in electronic form. To be trustworthy, it must
be authenticated.

\textbf{Authentication} proves that a message or document is genuine and
comes from the claimed sender. A \textbf{digital signature} provides
authenticity, integrity, and non-repudiation. \textbf{Encryption}, by
contrast, provides confidentiality: it keeps the content secret from
unauthorised readers. \textbf{Aadhaar} is used as a digital identity for
authentication of a person.

\begin{definitionbox}{Digital signature}
A mechanism used to authenticate an electronic record. It provides
authenticity, integrity, and non-repudiation, and is distinct from
encryption, which provides confidentiality.
\end{definitionbox}

\begin{example}
An item states, ``A digital signature encrypts the document so that only the
recipient can read it.'' This confuses two ideas. Encrypting for secrecy is
\textbf{encryption}, which provides confidentiality. A \textbf{digital
signature} proves who sent the document and that it was not altered; it
provides authenticity and integrity. The wrong option has used the
signature's name for encryption's job.
\end{example}

\section{Exam retrieval and common confusions}
Six distinctions prevent most errors on this topic.
\begin{enumerate}
  \item e-NAM = agricultural marketing and trading; it is \textbf{not}
    general e-commerce.
  \item GeM = G2B procurement, \textbf{not} a G2C service. e-Office and PFMS
    are G2G.
  \item NeGP (plan, 2006), e-Kranti (pillar / NeGP 2.0, 2015), and NeGD
    (division, 2009) are three different things; do not swap them.
  \item Digital signature = authenticity and integrity; encryption =
    confidentiality. They are \textbf{not} the same (TRAP-15).
  \item The citizen portal of the Government of India is india.gov.in, not
    the Digital India or NeGD site.
  \item Digital India is built on three vision areas and nine pillars;
    e-Kranti is Pillar 5, the service-delivery pillar.
\end{enumerate}

\begin{example}
An item asks which body supports MeitY in the programme management of
e-Governance projects. Trace the three names: NeGP is a plan, e-Kranti is a
pillar, and only NeGD is an organisation. The answer is NeGD.
\end{example}

\subsection*{Exercises}
\begin{enumerate}[label=\arabic*.]
  \item Expand e-NAM and state its purpose in one line.
  \item Distinguish G2C, G2G, G2B, and G2E, giving one example of each.
  \item State the launch date and vision of the Digital India programme.
  \item Explain the difference between NeGP, e-Kranti, and NeGD.
  \item Define authentication and explain how a digital signature differs
    from encryption.
\end{enumerate}

\subsection*{Solutions}
\begingroup\small
\begin{enumerate}[label=\arabic*.]
  \item e-NAM is the National Agriculture Market: a pan-India electronic
    trading portal that networks the existing APMC mandis to create a
    unified national market for agricultural commodities.
  \item G2C (Government to Citizen) serves the public, e.g.\ UMANG; G2G
    (Government to Government) serves government bodies, e.g.\ e-Office;
    G2B (Government to Business) serves businesses, e.g.\ GeM; G2E
    (Government to Employee) serves government employees, e.g.\ an employee
    service book.
  \item Digital India was launched on 1 July 2015, with the vision to
    transform India into a digitally empowered society and knowledge
    economy.
  \item NeGP is the National e-Governance Plan, the umbrella plan approved
    in May 2006 with 27 MMPs. e-Kranti is Pillar 5 of Digital India and the
    upgraded NeGP (NeGP 2.0). NeGD is the National e-Governance Division, an
    organisation created in 2009 under MeitY to support e-Governance
    projects.
  \item Authentication proves that a message or document is genuine and
    comes from the claimed sender. A digital signature provides authenticity,
    integrity, and non-repudiation; encryption provides confidentiality by
    keeping the content secret. They are different functions.
\end{enumerate}
\endgroup
\end{document}
